\chapter{Summary}
\label{CH:SUM}

Este texto apresentou os conceitos fundamentais de sistemas operacionais por meio de um estudo detalhado, linha por linha, de um único sistema: o xv6. Algumas linhas de código representam a essência de conceitos fundamentais de sistemas operacionais—como troca de contexto, a fronteira entre o espaço do usuário e o do kernel, e mecanismos de bloqueio—e são cruciais para entender como um sistema operacional funciona. Outras linhas servem como exemplos ilustrativos de como implementar certas ideias, e poderiam ser desenvolvidas de maneiras diferentes—por exemplo, com algoritmos de escalonamento mais eficientes, estruturas de dados em disco mais eficazes, ou mecanismos de registro (logging) aprimorados para permitir transações concorrentes.

Todos os conceitos foram explorados sob a perspectiva de uma interface de chamadas de sistema específica e amplamente consagrada: a interface Unix. Embora a implementação tenha sido específica ao xv6, as ideias fundamentais se aplicam amplamente ao projeto e à compreensão de diversos sistemas operacionais modernos.
